% Introduction

Vessel diameter is an important biomarker that can provide valuable information regarding local vascular health. Thickening of the arteriolar wall has been shown to narrow the vessel lumen in proliferative diabetic retinopathy~\cite{uenoAssociationChangesRetinal2021}. Retinal arteriolar narrowing has been shown to precede and predict the onset of systemic hypertension~\cite{ikramRetinalVesselDiameters2006}. In glaucoma suspects, reduced vessel diameter has been shown to track both retinal ganglion cell dysfunction and visual field loss~\cite{tirsiRetinalVesselDiameter2025}. Diameter changes are also informative beyond the eye, where a meta-analysis of observational studies found smaller retinal arteriolar caliber to be associated with coronary heart disease~\cite{liangRetinalVascularCaliber2022}. Different diseases have been shown to preferentially affect different plexuses, with glaucoma altering the superficial vascular complex~\cite{takusagawaProjectionResolvedOptical2017} while diabetic retinopathy~\cite{hwangVisualization3Distinct2016} and retinal vein occlusion~\cite{coscasOpticalCoherenceTomography2016} affect the deeper plexuses more severely. Combined with measurements of blood flow, diameter can also be used to investigate local blood pressure and vessel elasticity~\cite{kannenkerilRetinalVascularResistance2018,mynardMeasurementAnalysisInterpretation2020,seabraSilicoBloodPressure2022}. For these reasons, it is imperative that accurate measurements of vessel diameter be used for these extrapolations or when modeling behaviours of flow in digital twins~\cite{hernandezLinkingVascularStructure2024,hernandezMicrovascularRetinalDigital2026}. Optical imaging techniques are often used to measure vessel diameter due to their high resolution, but accurate measurement of vessel size is more difficult than it first appears.

Optical coherence tomography (OCT) is a non-invasive optical imaging method, which uses interferometry to measure the transit time of light that is backscattered from a sample, thereby producing depth-resolved images of tissue~\cite{huangOpticalCoherenceTomography1991,fercherVivoOpticalCoherence1993}. Optical coherence tomography angiography (OCTA) is an extension of OCT that visualizes vascular perfusion using the motion contrast from red blood cells (RBCs)~\cite{makitaOpticalCoherenceAngiography2006,srinivasanRapidVolumetricAngiography2010,taoSinglepassVolumetricBidirectional2008,wangThreeDimensionalOptical2007}. Over the years, OCTA has demonstrated widespread acceptance and use for clinical evaluations of vascular features in relation to disease~\cite{javedOpticalCoherenceTomography2023}. OCT angiograms can be used to evaluate characteristics on the scale of vascular layers by measuring vascular density or non-perfusion zones~\cite{tanApproachesQuantifyOptical2020,hwangVisualization3Distinct2016,coscasOpticalCoherenceTomography2016,takusagawaProjectionResolvedOptical2017,zabelComparisonRetinalMicrovasculature2019,tanApplicationOpticalCoherence2021}, or it can be used to guage vessel size~\cite{yaoComparisonRetinalVessel2021} or tortuosity~\cite{leeQUANTIFICATIONRETINALVESSEL2018}.

Analysis of OCTA datasets is typically achieved by performing a maximum intensity projection of a 3D vascular network to compress it into 2D before binarising the resulting angiogram using a global or local binarisation method. From this binarised angiogram, various vascular metrics can then be quantified. Importantly, there are many methods for performing this binarisation step and the choice of method can significantly affect the reported vascular metrics, with no clear consensus for a one size fits all approach~\cite{mehtaImpactBinarizationThresholding2019,rabioloComparisonMethodsQuantify2018,borrelliOPTICALCOHERENCETOMOGRAPHY2021,freedmanImpactImageProcessing2022,arrigoImpactDifferentThresholds2021}. Similarly, there are vascular enhancement algorithms to highlight structures that resemble vessels before quantification, but these methods also risk affecting the apparent size of the vessels in the network~\cite{freedmanImpactImageProcessing2022,hongIntrasessionRepeatabilityQuantitative2020,hongEffectVesselEnhancement2020}.

Furthermore, OCT is prone to certain artifacts which specifically affect visualization of the vasculature~\cite{cimallaShearFlowinducedOptical2011,zhuVisibilityMicrovesselsOptical2020,bernucciInvestigationArtifactsRetinal2018}. When viewed in cross-section, large vessels often demonstrate an hourglass appearance with stronger scattering in the center than at the sides. This appearance is due to orientation effects of RBCs, which align with the blood flow, causing a reduced scattering cross-section at the sides of the vessels~\cite{cimallaShearFlowinducedOptical2011}. We hypothesise that these artifacts may lead to underestimations of measurements of vessel diameter. To compensate for these artifacts, OCT contrast agents such as Intralipid have previously been used to increase intravascular intensity signals, particularly in regions where RBC scattering is diminished~\cite{panUltrasensitiveDetection3D2014,merkleDynamicContrastOptical2016,merkleLaminarMicrovascularTransit2016,zhuVisibilityMicrovesselsOptical2020,bernucciInvestigationArtifactsRetinal2018}.  

In this study we examine different methods of measuring vessel diameter from in vivo OCT scans of mouse retinal vasculature, and we use a contrast agent to obtain ground truth vessel diameter measurements for comparison. In addition to traditional binarisation-based methods, we also employ model-based methods for quantifying vessel diameter. Traditional Gaussian models have been employed for vessel diameter measurements in techniques such as fluorescein angiography~\cite{zhouDetectionQuantificationRetinopathy1994} and fundus photography~\cite{lowellMeasurementRetinalVessel2004} (based on other optical imaging methods) and OCT~\cite{leahyMapping3DConnectivity2015}, but we expand this by also investigating the generalized Gaussian model. As part of this study, we measure how much diameter is underestimated from enface projections and evaluate how reliably it can be compensated through calibration. Flicker stimulation is well established to evoke vasodilation in the retina~\cite{garhoferDiffuseLuminanceFlicker2004,sharifizadFactorsDeterminingFlickerinduced2016,michelsonFlickeringLightIncreases2002,aschingerEffectofDiffuse2017,kornfieldRegulationBloodFlow2014,shihQuantitativeRetinalChoroidal2013,sonOpticalCoherenceTomography2016}. After identifying the best methods for quantifying diameter and compensating for underestimations, we finally evaluated a test scenario using flicker stimulation to induce a measurable change in diameter to evaluate the performance of different methods.

%\section{Introduction}

% General structure of intro:
% Why diameter is an important measurment
% Brief OCTA Introduction
% Difficulty of OCTA quantification (different methods yield different results)
% Artefact contributions and diameter underestimation
% Proposed solution that we investigate here and evaluation of different methods in a test scenario (vessel dilation)

% https://www.nature.com/articles/s41598-021-84067-2
% https://qims.amegroups.org/article/view/51210/html
% https://jamanetwork.com/journals/jamaophthalmology/fullarticle/415121
% https://www.sciencedirect.com/science/article/pii/S0307904X23002913
% https://www.pnas.org/doi/10.1073/pnas.2521872123

%In addition to thresholding-based variability, we hypothesize that OCTA metrics are also affected by OCT artifacts caused by the non-uniform backscattering profiles of red blood cells (RBCs).  Because RBCs have a distinct bi-concave disc shape, they more effectively backscatter light when aligned perpendicular to the beam, rather than parallel to it.  It has already been observed that in large vessels, a so-called hourglass artifact appears when viewed cross-sectionally.  This is caused by the alignment of the RBCs with the vessel walls while under laminar flow, where the RBCs are perpendicular to the beam in the center of the vessels and parallel to the beam at the edges of the vessels.  This artifact artificially suppresses the OCT and OCTA signals at the edges of the vessels, making them appear more narrow when projected into an en-face view, which is the preferred way of displaying angiographic data.  Here we use contrast enhanced OCT (CE-OCT) to investigate how significant these effects are and how they can be corrected.

%Optical coherence tomography (OCT) is a non-invasive optical imaging method, which uses interferometry to measure the transit time of light that is backscattered from a sample.  OCT has found widespread clinical use in ophthalmology with XXX number of systems and XXX OCT scans performed annually around the world.  OCT angiography (OCTA) is a functional extension of OCT which highlights motion between repeated OCT scans to generate images of vascular networks.  Typically, OCTA signals are first calculated from the OCT intensity and/or phase information, segmented in 3D manually, algorithmically, or using a deep learning approach, and finally compressed into a 2D image using a maximum intensity projection.  These 2D vascular maps can then be binarised by using a global or adaptive threshold, and finally vascular metrics such as vessel density and diameter can be computed.  Of these processing stages, the thresholding step is perhaps the most important.  There are dozens of different thresholding methods that have been proposed over the years and each of these may produce different results in the binary map, which would in turn influence the quantification of vascular metrics. In this work, we examine an alternate approach for extracting quantitative vessel diameter by fitting a model to the OCTA data before binarisation to avoid thresholding-based effects.


%Intralipid has been shown to be a highly scattering contrast agent that fills the vascular lumen~completely, thereby compensating for the artefact caused by the diminished backscattering at vessel peripheries caused by RBC orientation effects~\cite{merkleHighresolutionDepthresolvedVascular2021,panUltrasensitiveDetection3D2014,cimallaShearFlowinducedOptical2011}.

%The physiological response to flicker stimulation produces vasodilation, which has been well established in retinal physiology~\cite{garhoferDiffuseLuminanceFlicker2004,sharifizadFactorsDeterminingFlickerinduced2016}.